% ====================================================================
% 서지 — 29종 (tier 등급·과인용 가드 주석 포함, 이전본 계승; v1.0.20 전수 실재 검증 — results/V1020_REFERENCE_LEDGER.md)
% ====================================================================
\begin{thebibliography}{99}
\bibitem{dahn1991} J. R. Dahn, ``Phase diagram of Li$_x$C$_6$,'' \emph{Phys. Rev. B} \textbf{44}, 9170 (1991). DOI: 10.1103/PhysRevB.44.9170.
\bibitem{ohzuku1993} T. Ohzuku, Y. Iwakoshi, K. Sawai, ``Formation of Lithium-Graphite Intercalation Compounds in Nonaqueous Electrolytes and Their Application as a Negative Electrode for a Lithium Ion (Shuttlecock) Cell,'' \emph{J. Electrochem. Soc.} \textbf{140}, 2490 (1993). DOI: 10.1149/1.2220849.
\bibitem{bazant2013} M. Z. Bazant, ``Theory of Chemical Kinetics and Charge Transfer based on Nonequilibrium Thermodynamics,'' \emph{Acc. Chem. Res.} \textbf{46}, 1144 (2013). DOI: 10.1021/ar300145c.
\bibitem{eyring1935} H. Eyring, ``The Activated Complex in Chemical Reactions,'' \emph{J. Chem. Phys.} \textbf{3}, 107 (1935). DOI: 10.1063/1.1749604.
\bibitem{hill1960} T. L. Hill, \emph{An Introduction to Statistical Thermodynamics}, Addison--Wesley (1960); Dover reprint (1986) --- Ch.~7 (국소화 흡착의 대정준 유도, $q(T)$ 와 Langmuir 등온선).
\bibitem{fowler1939} R. H. Fowler and E. A. Guggenheim, \emph{Statistical Thermodynamics}, Cambridge Univ. Press (1939) --- 국소화 단분자층 흡착의 통계역학 유도.
\bibitem{mcquarrie1976} D. A. McQuarrie, \emph{Statistical Mechanics}, Harper \& Row (1976); Univ. Science Books reprint (2000) --- 독립 부분계 분배함수의 인수분해와 흡착 등온선.
\bibitem{ashcroftmermin1976} N. W. Ashcroft and N. D. Mermin, \emph{Solid State Physics}, Holt, Rinehart and Winston (1976) --- 전자 기체의 Fermi--Dirac 통계와 Sommerfeld 전개(Ch.~2·Appendix C, 장 수준 인용).
\bibitem{mckinnon1983} W. R. McKinnon and R. R. Haering, ``Physical Mechanisms of Intercalation,'' in \emph{Modern Aspects of Electrochemistry}, No.~15 (R. E. White, J. O'M. Bockris, B. E. Conway, eds.), Plenum Press (1983), pp.~235--304 --- 격자기체 삽입 등온선의 원전.
\bibitem{dreyer2010} W. Dreyer \emph{et al.}, ``The thermodynamic origin of hysteresis in insertion batteries,'' \emph{Nature Materials} \textbf{9}, 448 (2010). DOI: 10.1038/nmat2730.
\bibitem{bloom2005} I. Bloom \emph{et al.}, ``Differential voltage analyses of high-power, lithium-ion cells,'' \emph{J. Power Sources} \textbf{139}, 295 (2005). DOI: 10.1016/j.jpowsour.2004.07.021.
\bibitem{dubarry2012} M. Dubarry, C. Truchot, B. Y. Liaw, ``Synthesize battery degradation modes via a diagnostic and prognostic model,'' \emph{J. Power Sources} \textbf{219}, 204 (2012). DOI: 10.1016/j.jpowsour.2012.07.016.
% --- LCO 양극 참고문헌 ---
\bibitem{reimers1992} J. N. Reimers and J. R. Dahn, ``Electrochemical and in situ X-ray diffraction studies of lithium intercalation in Li$_x$CoO$_2$,'' \emph{J. Electrochem. Soc.} \textbf{139}, 2091--2097 (1992). DOI: 10.1149/1.2221184.
\bibitem{menetrier1999} M. M\'en\'etrier, I. Saadoune, S. Levasseur, and C. Delmas, ``The insulator--metal transition upon lithium deintercalation from LiCoO$_2$: electronic properties and $^7$Li NMR study,'' \emph{J. Mater. Chem.} \textbf{9}, 1135--1140 (1999). DOI: 10.1039/a900016j.
\bibitem{motohashi2009} T. Motohashi, T. Ono, Y. Sugimoto, Y. Masubuchi, S. Kikkawa, R. Kanno, M. Karppinen, and H. Yamauchi, ``Electronic phase diagram of the layered cobalt oxide system Li$_x$CoO$_2$ ($0\le x\le1$),'' \emph{Phys. Rev. B} \textbf{80}, 165114 (2009). DOI: 10.1103/PhysRevB.80.165114; arXiv:0909.3556.
\bibitem{xia2007} H. Xia, L. Lu, Y. S. Meng, and G. Ceder, ``Phase transitions and high-voltage electrochemical behavior of LiCoO$_2$ thin films grown by pulsed laser deposition,'' \emph{J. Electrochem. Soc.} \textbf{154}(4), A337--A342 (2007). DOI: 10.1149/1.2509021.
\bibitem{reynier2004} Y. Reynier, J. Graetz, T. Swan-Wood, P. Rez, R. Yazami, and B. Fultz, ``Entropy of Li intercalation in Li$_x$CoO$_2$,'' \emph{Phys. Rev. B} \textbf{70}, 174304 (2004). DOI: 10.1103/PhysRevB.70.174304. [동 그룹: \emph{J. Electrochem. Soc.} \textbf{151}, A422 (2004).]
\bibitem{swiderska2019} A. {\'S}widerska-Mocek, E. Rudnicka, and A. Lewandowski, ``Temperature coefficients of Li-ion battery single electrode potentials and related entropy changes---revisited,'' \emph{Phys. Chem. Chem. Phys.} \textbf{21}, 2115 (2019). DOI: 10.1039/c8cp06638h. [LiCoO$_2$ 단전극 $\partial\phi/\partial T\approx+0.83$ mV/K, tier B.]
\bibitem{msmr2024} A. Paul, K. Wolfe, M. W. Verbrugge, B. J. Koch, J. S. Lowe, J. Trembly, J. A. Staser, and T. R. Garrick, ``Quantifying the Temperature Dependence of the Multi-Species, Multi-Reaction Model. Part~1: Parameterization for a Meso-Carbon Micro-Bead Graphite,'' \emph{ECS Advances} \textbf{3}, 042501 (2024). DOI: 10.1149/2754-2734/ad7d1c. [Part~II: \emph{J. Electrochem. Soc.} \textbf{171}(10), 103505 (2024), DOI: 10.1149/1945-7111/ad70d9.]
\bibitem{ml2024} M. Faghih Shojaei, J. Holber, S. Das, G. H. Teichert, T. Mueller, L. Hung, V. Gavini, and K. Garikipati, ``Bridging scales with Machine Learning: From first principles statistical mechanics to continuum phase field computations to study order--disorder transitions in Li$_x$CoO$_2$,'' \emph{J. Mech. Phys. Solids} \textbf{190}, 105726 (2024). DOI: 10.1016/j.jmps.2024.105726; arXiv:2302.08991.
% --- broadening 절 참고문헌 ---
\bibitem{leviaurbach1999} M. D. Levi and D. Aurbach, ``Frumkin intercalation isotherm --- a tool for the description of lithium insertion into host materials: a review,'' \emph{Electrochim. Acta} \textbf{45}, 167--185 (1999). DOI: 10.1016/S0013-4686(99)00202-9. [내재 평형 폭$\cdot$Frumkin 등온선.]
\bibitem{rsc2021} H. Onuma, K. Kubota, S. Muratsubaki, W. Ota, M. Shishkin, H. Sato, K. Yamashita, S. Yasuno, and S. Komaba, ``Phase evolution of electrochemically potassium intercalated graphite,'' \emph{J. Mater. Chem. A} \textbf{9}, 11187--11200 (2021). DOI: 10.1039/D0TA12607A. [저자$\cdot$제목$\cdot$권쪽은 Crossref 로 확정. 흑연 staging 전이별 1차/연속 구분(dilute$\cdot$4L--3L $=$ solid-solution plateau 없음, 나머지 $=$ two-phase plateau)은 이 문헌이 \emph{비교}로 정리한 흑연 삽입 상전이 순서에서 읽음 --- K-graphite 가 주제이나 Li-graphite staging 순서를 대조 제시(tier B, 비교 인용).]
\bibitem{fly2020} A. Fly and R. Chen, ``Rate dependency of incremental capacity analysis (d$Q$/d$V$) as a diagnostic tool for lithium-ion batteries,'' \emph{J. Energy Storage} \textbf{29}, 101329 (2020). DOI: 10.1016/j.est.2020.101329. [C-rate$\uparrow$ 일수록 peak 고전위 이동$\cdot$둔화 --- 동역학 비대칭 꼬리. 저자$\cdot$제목$\cdot$DOI 연도는 Crossref 로 확정(저자 2인 A.~Fly, R.~Chen --- 권 \textbf{29}, 아티클 101329).]
\bibitem{dahn1995} J. R. Dahn, T. Zheng, Y. Liu, and J. S. Xue, ``Mechanisms for Lithium Insertion in Carbonaceous Materials,'' \emph{Science} \textbf{270}(5236), 590--593 (1995). DOI: 10.1126/science.270.5236.590. [탄소 음극의 구조 무질서$\cdot$흑연화도와 리튬 삽입의 \emph{intra-particle} 용량 한계 통계($x_{\max}=1-P$ site-blocking 무질서 모델). \emph{용량-한계 예시}로 인용 --- 비-크기 구조 무질서가 실재함은 보이나, inter-particle 의 apparent-U($\eta$) 분포 형상을 직접 정량하는 1차 근거는 아님(tier C 근거미발견).]
\bibitem{park2021} J. H. Park, H. Yoon, Y. Cho, and C.-Y. Yoo, ``Investigation of Lithium Ion Diffusion of Graphite Anode by the Galvanostatic Intermittent Titration Technique,'' \emph{Materials} \textbf{14}(16), 4683 (2021). DOI: 10.3390/ma14164683. [흑연 하프셀 GITT --- 준평형 OCP 가 입자 크기$\cdot$전극 형상 무관 staging 전위 상수임을 보고; $D_\mathrm{Li}\approx10^{-11}$--$4\times10^{-10}$ cm$^2$/s. 서지는 Crossref 재확인.]
\bibitem{reynier2003} Y. Reynier, R. Yazami, and B. Fultz, ``The entropy and enthalpy of lithium intercalation into graphite,'' \emph{J. Power Sources} \textbf{119--121}, 850--855 (2003). DOI: 10.1016/S0378-7753(03)00285-4. [흑연 삽입 엔트로피$\cdot$엔탈피 실측 --- 부호$\cdot$스케일 anchor(tier B), 전이별 정밀값 아님.]
\bibitem{persson2010} K. Persson, V. A. Sethuraman, L. J. Hardwick, Y. Hinuma, Y. S. Meng, A. van der Ven, V. Srinivasan, R. Kostecki, and G. Ceder, ``Lithium Diffusion in Graphitic Carbon,'' \emph{J. Phys. Chem. Lett.} \textbf{1}, 1176--1180 (2010). DOI: 10.1021/jz100188d. [제일원리$+$실측 Li-흑연 확산 --- 이동 장벽 스케일(tier C).]
\bibitem{persson2010b} K. Persson, Y. Hinuma, Y. S. Meng, A. Van der Ven, and G. Ceder, ``Thermodynamic and kinetic properties of the Li-graphite system from first-principles calculations,'' \emph{Phys. Rev. B} \textbf{82}, 125416 (2010). DOI: 10.1103/PhysRevB.82.125416. [Li-흑연 열역학$\cdot$동역학 제일원리 --- 조성 의존 경향의 근거(tier C).]
\bibitem{cogswell2012} D. A. Cogswell and M. Z. Bazant, ``Coherency Strain and the Kinetics of Phase Separation in LiFePO$_4$ Nanoparticles,'' \emph{ACS Nano} \textbf{6}, 2215--2225 (2012). DOI: 10.1021/nn204177u. [상분리 입자 전극의 계면 에너지 논의 참조 --- 본 문건 $\gamma$ 수치의 출처 아님(스케일 상정은 tier C).]
\end{thebibliography}

% 자산: [B2-154] [B2-155] [B2-156] [B2-157] [B2-158] [B2-159] [B2-160] [B2-161] [B2-162] [B2-163] [B2-164] [B2-165] [B2-166] [B2-167] [B2-168] [B2-169] [B2-170] [B2-171] [B2-172] [B2-173] [B2-174] [B2-175] [B2-176] [B2-177]
